\subsection{トピックと参考文献の対応}
この表は、SWEBOK が示す「トピック対参考資料の行列」を、学習用に簡略化したものである。詳しく学ぶときは、各トピックに対応する文献、規格、公開データベースを確認するとよい。

\par\smallskip\noindent\refstepcounter{table}\label{tab:ch13-06}表 \thetable: トピックと参考文献の対応の整理\par\smallskip
表 \thetable は、トピックと参考文献の対応の整理を比較・確認しやすい形で整理したものである。\par\smallskip
\begin{longtable}{>{\raggedright\arraybackslash}p{0.25\linewidth}>{\raggedright\arraybackslash}p{0.7\linewidth}}
\toprule
トピック & 主な参考資料 \\
\midrule
\endhead
1. 基礎 & ISO/IEC 25010, ISO/IEC 27000, ISO/IEC 27032, Easttom, Diogenes and Ozkaya \\
2. 管理と組織 & Allen et al.; Bishop; ISO/IEC 27001; ISO/IEC 21827; Agile Application Security \\
3. セキュリティ工学とプロセス & Allen et al.; Howard and Lipner; DoD Enterprise DevSecOps; ISO/IEC 15408 \\
4.1 セキュリティ要求 & Allen et al.; McGraw; Bishop; Firesmith \\
4.2 セキュリティ設計 & Allen et al.; Swiderski and Snyder; Bishop; security design references \\
4.3 セキュリティパターン & Fernandez-Buglioni; Nagappan et al.; Schumacher et al. \\
4.4 セキュリティを考慮した構築 & Seacord; CERT secure coding references; Allen et al.; Bishop \\
4.5 セキュリティテスト & Allen et al.; McGraw; Bishop; NIST SP800-115; PCI DSS \\
4.6 脆弱性管理 & CVE; CWE; CAPEC; CVSS; Bishop \\
5. ツール & Allen et al.; Erickson; McGraw; Dotson \\
6. 領域別セキュリティ & Dotson; IEEE IoT Security Best Practices; IEC IoT 2020; Chio and Freeman \\
\bottomrule
\end{longtable}
